% Test-case specification TC-0001 (CTest test TC-0001-library).
\documentclass[letterpaper]{article}

\usepackage[letterpaper,margin=1in]{geometry}
\emergencystretch=3em

% Compile from the repository root so this path resolves to
% the pinned WSP submodule.
\input{wsp/testing/test-case-library.tex}

\begin{document}

\def\TCID{TC-0001}
\def\TCTitle{%
  Export directory, loader and engine-entry diagnostic}
\def\TCPurpose{%
  Verify that the built DLL has the export directory of the reference
  runtime (ordinal base 100, 1,911 slots, 635 named exports at their
  reference ordinals, the other slots empty), that it loads with names and
  ordinals resolving, and that the unsupported p-code entry stops the
  process with its diagnostic.}
\def\TCPriority{\TCPriorityCritical}

\def\TCRequirementRef{%
  REQ-0001, REQ-0002, REQ-0007}
\def\TCDesignRef{%
  Not applicable.}
\def\TCFeatureRef{%
  \texttt{tests/\allowbreak{}vbrun\_\allowbreak{}tests.\allowbreak{}cpp},
  \texttt{tests/\allowbreak{}run\_\allowbreak{}vbrun\_\allowbreak{}tests.\allowbreak{}cmake},
  \texttt{src/\allowbreak{}wfcvbrun.\allowbreak{}def},
  \texttt{src/\allowbreak{}vbrun\_\allowbreak{}exports.\allowbreak{}inc}}
\def\TCTestTechnique{%
  Back-to-back comparison with the retained reference inventory; boundary
  values at the first and last ordinal and at a hole.}

\def\TCPreconditions{%
  The library is built. The test is registered only when the retained
  inventory
  \texttt{evidence/\allowbreak{}reference/\allowbreak{}msvbvm60-6.\allowbreak{}0-runtime.\allowbreak{}json}
  exists (a WFC checkout, or
  \texttt{WFC\_\allowbreak{}VBRUN\_\allowbreak{}EVIDENCE} names it); the
  standalone build does not register it.}
\def\TCEnvironment{%
  Windows with the toolchain of
  \texttt{scripts/\allowbreak{}Invoke-Build.\allowbreak{}ps1}; the local
  gate builds and runs x64 Debug, CI builds and runs the other
  architectures. The library is built from this repository in
  \texttt{out/\allowbreak{}build}.}
\def\TCAssumptions{%
  Not applicable.}

\def\TCInputData{%
  The inventory, converted by
  \texttt{run\_\allowbreak{}vbrun\_\allowbreak{}tests.\allowbreak{}cmake}
  into \texttt{exports.\allowbreak{}csv} (one \texttt{ordinal,name} line per
  present export); ordinals 100 (\texttt{ThunRTMain}), 729 (a hole) and 2010
  (the last slot); the export \texttt{ProcCallEngine}.}
\def\TCInitialState{%
  A fresh work folder \texttt{vbrun-test} in the build tree.}

\def\TCProcedure{%
  \begin{enumerate}
    \item CTest runs the test program.
    \item CTest runs \texttt{cmake -P
      tests/\allowbreak{}run\_\allowbreak{}vbrun\_\allowbreak{}tests.\allowbreak{}cmake}
      with the test program, the DLL, the inventory and the work folder.
    \item \texttt{vbrun\_\allowbreak{}tests layout} parses the PE export
      directory of the DLL and compares ordinal base, table size, names,
      ordinals and holes with \texttt{exports.\allowbreak{}csv}.
    \item \texttt{vbrun\_\allowbreak{}tests load} loads the DLL, resolves
      \texttt{ThunRTMain} by name and by ordinal 100, resolves ordinal 2010,
      and looks up ordinal 729.
    \item \texttt{vbrun\_\allowbreak{}tests call} calls the export
      \texttt{ProcCallEngine}.
    \item CTest records the exit code and, on failure, the program's output
      (\texttt{--output-on-failure}).
  \end{enumerate}}
\def\TCExpectedResult{%
  Step 2 prints \texttt{layout ok} with 635 named exports and exits 0. Step
  3 prints \texttt{load ok} and exits 0: the name and ordinal 100 give the
  same address, ordinal 2010 resolves and ordinal 729 does not. Step 4 stops
  the process with exit code 22082 (0x5642) and the output contains
  \texttt{p-code (ProcCallEngine) is not supported}. Any other result fails
  the test.}
\def\TCPostConditions{%
  The work folder holds \texttt{exports.\allowbreak{}csv}; the DLL is
  unloaded when each process ends.}

\TCTemplate

\end{document}
